\subsection{Positive linear forms on a Riesz space}

We recall the following definition (TVS, II, §2, No.~5):

DEFINITION 1. --- Given an ordered vector space E, a linear form L on E is said to be positive if \(L(x) \geqslant 0\) for every \(x \geqslant 0\) in E.

Since \(L(y) - L(x) = L(y - x)\) , it is the same to say that the relation \(x \leqslant y\) implies \(L(x) \leqslant L(y)\) , or again that \(L\) is an increasing function on \(\mathbf{E}\) .

Examples. --- 1) Let A be any set, E a linear subspace of the space \(R^{A}\) of all real-valued functions defined on A. For every element \(a \in A\) , the mapping \(x \mapsto x(a)\) is a positive linear form on E.

\begin{enumerate}
\def\labelenumi{\arabic{enumi})}
\setcounter{enumi}{1}
\item
  Let \(\mathrm{I} = [a, b]\) be a compact interval of \(\mathbf{R}\) , \(\mathrm{E}\) the Riesz space formed by the regulated real-valued functions on \(\mathrm{I}\) (FRV, II, §1, No.~3); the mapping \(x \mapsto \int_{a}^{b} x(t) dt\) is a positive linear form on \(\mathrm{E}\) .
\item
  Let F be any set, U an ultrafilter on F (GT, I, §6, No.~4), E the Riesz space \(\mathcal{B}(F)\) of bounded real-valued functions on F. For every \(x \in E\) , \(\lim_{\mathfrak{U}} x(t)\) exists, because \(x(\mathfrak{U})\) is a base of an ultrafilter on the relatively compact set \(x(F)\) , hence is convergent. Moreover, if \(x \geqslant 0\) then \(\lim_{\mathfrak{U}} x(t) \geqslant 0\) by the principle of extension of inequalities; the mapping \(x \mapsto \lim_{\mathfrak{U}} x\) is thus a positive linear form on E. If U is taken to be the ultrafilter formed by the sets containing an element \(a \in F\) , one recovers the positive linear form \(x \mapsto x(a)\) (Example 1).
\end{enumerate}

PROPOSITION 1. --- Let E be an ordered vector space, L a mapping of E into R such that \(L(x+y)=L(x)+L(y)\) and such that the relation \(x\geqslant0\) implies \(L(x)\geqslant0\) ; then \(L(\lambda x)=\lambda L(x)\) for every scalar \(\lambda\) and every \(x\geqslant0\) .

Since \(L(-x) = -L(x)\) (L being a representation of the additive group E in R), we can restrict ourselves to the case that \(\lambda \geqslant 0\) . For every integer \(n \geqslant 0\) , we have \(L(nx) = nL(x)\) , whence \(L((1/n)x) = (1/n)L(x)\) and consequently \(L(rx) = rL(x)\) for every rational number \(r \geqslant 0\) . On the other hand, L is increasing in E; if r and \(r'\) are rational numbers such that \(r \leqslant \lambda \leqslant r'\) , it follows that \(rL(x) \leqslant L(\lambda x) \leqslant r'L(x)\) ; since \(rL(x)\) and \(r'L(x)\) differ from \(\lambda L(x)\) as little as we like, we have \(L(\lambda x) = \lambda L(x)\) .

PROPOSITION 2. --- Let E be a real vector space, C a convex cone with vertex 0 in E such that E = C - C, and \(x \mapsto M(x)\) a mapping of C into R such that \(M(\lambda x + \mu y) = \lambda M(x) + \mu M(y)\) for all \(x \in C\) , \(y \in C\) , \(\lambda \geqslant 0\) , \(\mu \geqslant 0\) . Then, there exists one and only one linear form L that extends M to E.

By hypothesis, every \(z \in E\) may be written \(z = y - x\) , where x, y belong to C; moreover, if also \(z = y' - x'\) with \(x' \in C\) , \(y' \in C\) , then

\(M(y) - M(x) = M(y') - M(x')\) ; for, from the relation \(y - x = y' - x'\) we infer \(y + x' = x + y'\) , consequently \(M(y) + M(x') = M(x) + M(y')\) . Let us denote by \(L(z)\) the common value of \(M(y) - M(x)\) for all expressions of \(z\) as the difference \(y - x\) of two elements of C; one verifies immediately that \(L\) is a linear form on E extending \(M\) ; the uniqueness of \(L\) results from the fact that C generates the space E.

PROPOSITION 3. --- Let E be a directed ordered vector space, P the set of elements \(\geqslant 0\) of E, and \(x \mapsto M(x)\) a mapping of P into R, with values \(\geqslant 0\) , such that \(M(x + y) = M(x) + M(y)\) for all x, y in P. Then, there exists one and only one positive linear form L that extends M to E.

Since E = P - P, the same reasoning as in Prop. 2 proves, first, the existence and uniqueness of an additive mapping L of E into R that extends M. Proposition 1 then shows that \(L(\lambda x) = \lambda L(x)\) for all \(\lambda \geqslant 0\) and all \(x \in P\) , from which it is immediate that L is a linear form.

